% Transcription — fonds Alexandre Grothendieck, shelfmark 8, batch 1 (pages 1-20).
% Established from the facsimile archives/batches/8/batch-01.pdf, cut from a
% copy of 8.pdf fetched through the project's own relay
% (https://grothendieck-relay.onrender.com/source/8.pdf) rather than from
% grothendieck.umontpellier.fr directly; per the skill transcribe-grothendieck.
% The batch file carries no cover sheet: PDF page k is the archivists' page k,
% and their pencil numbers bottom-left agree throughout. The folder is eighty
% pages; this is the first of four batches (the others are transcribed
% separately by other passes).
%
% Pass: Opus 5.5 (claude-opus-5-5), 2026-09-26 — first pass, unchecked against the pages by a human.
%
% The folder sits in the inventory group "4-9" « Cristaux »; it has its own
% date in the catalogue, which \dating{} carries verbatim, with no group
% sentence.
%
% Pages skipped: 2, an upside-down typed leaf in English (p. 132, « §24.
% Explicit description of pseudo-rational double points », on two-dimensional
% normal local rings), with nothing in his hand; 18, a Bourbaki « Tribu n° 66 »
% (« Décisions sur l'intégration »), nothing in his hand; 20, a faded typed
% leaf « - 65 - » on holomorphic functions (Schwarz's lemma), nothing in his
% hand. All three are unrelated versos, skipped by the settled rule.
%
% Pages given as typescript, one French \note each with the hand-made
% interventions in full: 6, 8, 10, 14. They are four leaves of one French
% typescript, numbered in ink 35, 37.1, 32 (upside down) and 33, on a proper
% base-change argument for torsion sheaves (a theorem (5.1)(i)-(iii) for
% triples (g, f, F), Lemmes (8.3), (8.4), effacement of (R^q f'_*) g'^*, the
% Leray spectral sequence). They carry ink corrections and marginal remarks.
% The typescript is not identified here (no reference is looked up, per the
% skill), nor is the hand of the corrections asserted; the notes describe the
% text and give the interventions, without re-setting the typed prose.
%
% Pages transcribed from his hand: 1, 3-5, 7, 9, 11-13, 15-17, 19. Fast
% scratch work (« gribouillis ») in blue and black ink, mostly diagrams; the
% formulas carry, much of the connective prose does not and is left \ill{}.
% Page 16 is cancelled by two long pencil diagonals and is transcribed under
% them.
%
% His pagination: none on these leaves.
%
% What the batch is about. The Dieudonné crystal of a finite locally free
% group scheme G (or a truncated Barsotti-Tate group) killed by p, over a base
% S of characteristic p, and how to describe it by "perfect complexes" on the
% Zariski site: phi : S_cris -> S_zar with its section epsilon, phi epsilon =
% Frobenius; the crystal M = D^*(G) with F, V (FV = VF = 0); the crystalline
% complex Delta^* = Delta^*(M) of perfect amplitude in [0,1], H^0 = H^1 = M,
% with two exact triangles (I), (II) built from [M^(p) -F-> M] and
% [M -V-> M^(p)] and an identification phi^*(Delta^*) with Delta^*(p); on the
% Zariski side two perfect complexes l., l'. of amplitude [-1,0] with a
% pairing (l.^{G*})^v -> l.^G and a triangle (III) l.[-1] -> Delta^* -> l'.
% whose pull-back by phi^* is (I). Pages 5, 7: compatibility of these
% triangles with the (F,V)-structures. Page 9: the same programme for a
% crystal of finite presentation over S_cris/Z_p, extension to a
% p-divisorial thickening S'. Page 11: questions a)-d) (are these data stacks
% on S_zar? is the "restriction" functor faithful? is omega -> structure fully
% faithful? can one attach a finite locally free group scheme to
% (omega, F_omega, V_omega)?). Pages 12-13: the case of truncated BT groups
% killed by p: 0 -> omega -> M -> nu' -> 0, 0 -> n -> M -> t' -> 0,
% Ker F_M = Im V_M, hence n = omega, t' = nu'. Page 15: a plan 1°-4° (crystal
% of finite presentation, triangles, "F-V-triangles", "F-V-isom."), p. 17 its
% restatement for (omega, F_omega, V_omega) with the question whether the
% functor is fully faithful. Page 16 (cancelled): indeterminacy of data
% (M, T) under mu_2-torsors over the adèles, H^1 = prod Q_p^*/Q_p^{*2}, and
% rigidity. Page 19: descent of the filtration, D_1, D_0, t^*, omega, Ext^1
% (S_zar mod S_cris; ...) as the place of the arbitrariness, the diagram
% J -> B -> A with lambda |-> lambda^p, and the simplicial scheme X x X x X
% over S with the Frobenius F_S.
%
% Notation fixed once for the batch: \mathcal{M} for his script M (the
% crystal), M for the roman M (its value on S_zar, = epsilon^*(\mathcal{M}));
% \mathbb{D}^*(G) for his double-stroked D; \Delta^* for his triangle-shaped
% Delta; the Frobenius twist ^{(p)} as he writes it (sometimes ^p); \ell_\bullet
% for his « l. » (script l with a dot), \check\ell for the l with a hacek
% (dual); S_{\mathrm{cris}}, S_{\mathrm{zar}}; \varphi, \varepsilon as on
% p. 1.
%
% Hardest pages and readings to check first: the second triangle of p. 1 and
% its top vertex; the box on p. 1 (« Π : ν^(p) -> ν »); the margin of p. 3
% and the ringed words of p. 3 b); the margin of p. 11 (left \ill{});
% p. 16 throughout (the adelic reading « Q mod A » is doubtful); p. 19,
% right-hand column and the struck lines of its middle.

\documentclass[11pt,a4paper]{article}
% Le préambule commun à toutes les transcriptions du fonds Grothendieck.
%
% Il tient en une idée : **l'incertitude se compose, elle ne se résout pas.**
% Une transcription qui lisse un mot douteux a détruit la seule chose pour
% laquelle on la faisait — un lecteur doit pouvoir savoir, à chaque endroit,
% ce qui était sur le feuillet et ce qui a été ajouté. Les six macros qui
% suivent sont donc l'appareil critique, et non de la décoration.
%
% Les mêmes commandes sont reconnues par `scripts/render.mjs`, qui produit la
% vue de lecture du site. Une macro ajoutée ici doit l'être aussi là-bas,
% faute de quoi le rendu échouera bruyamment — ce qui est voulu : un rendu
% silencieusement faux est pire qu'un rendu absent.

\ProvidesPackage{grothendieck}[2026/08/08 Transcriptions du fonds Grothendieck]

\usepackage{amsmath,amssymb,amsthm}
\usepackage{mathrsfs}
\usepackage{stmaryrd}
% Les diagrammes commutatifs sont partout dans ce fonds : c'est la langue
% même de la géométrie algébrique de Grothendieck, et une transcription qui
% les rendrait en prose ne transcrirait rien.
\usepackage{tikz-cd}
% Français par défaut — c'est la langue des feuillets — mais l'anglais est
% chargé aussi : la traduction et le résumé sont des documents anglais, et la
% typographie française (espaces avant les deux-points) y serait une faute.
% Ces documents basculent par \selectlanguage{english} en tête de corps.
\usepackage[english,french]{babel}
\usepackage{fontspec}
\usepackage{geometry}
\geometry{a4paper,margin=25mm,marginparwidth=28mm}
\usepackage{marginnote}
\usepackage{xcolor}
% ulem, not soul. `soul`'s \st and \ul are built on a character-by-character
% scan that XeLaTeX's font machinery breaks: the document fails with an
% undefined control sequence rather than degrading. ulem does the same two jobs
% and is Unicode-engine safe. `normalem` stops it hijacking \emph, which the
% transcriptions use for Grothendieck's own emphasis.
\usepackage[normalem]{ulem}
\usepackage{hyperref}
\hypersetup{colorlinks=true,linkcolor=black,urlcolor=black,pdfborder={0 0 0}}

% --- Métadonnées du lot -----------------------------------------------------
% Reprises telles quelles par le script de rendu, qui les affiche en tête de
% la vue de lecture. Elles fixent ce que le fichier prétend couvrir : sans
% elles, une transcription flotte sans point d'attache dans un fonds de seize
% mille feuillets.

\newcommand{\folder}[1]{\gdef\grothFolder{#1}}
\newcommand{\batch}[1]{\gdef\grothBatch{#1}}
\newcommand{\leaves}[2]{\gdef\grothFirst{#1}\gdef\grothLast{#2}}
\newcommand{\foldertitle}[1]{\gdef\grothFoldertitle{#1}}
\gdef\grothFolder{?}\gdef\grothBatch{?}\gdef\grothFirst{?}\gdef\grothLast{?}\gdef\grothFoldertitle{}

% --- Le résumé --------------------------------------------------------------
%
% Il ouvre la lecture modernisée, sous le titre « Résumé », et il est composé
% en petit. Ce n'est pas un ornement : c'est le seul passage du document qui
% ne soit pas des mathématiques, et un lecteur doit pouvoir le reconnaître
% avant de l'avoir lu — puis le sauter s'il connaît déjà le sujet, ou s'y
% arrêter s'il ne le connaît pas. Le filet qui le suit marque la reprise du
% corps.

\newenvironment{resume}
  {\small\setlength{\parskip}{0.45em}}
  {\par\medskip\hrule\medskip}

% --- Mots-clés ---------------------------------------------------------------
%
% En anglais, en clôture du résumé de la lecture modernisée : le vocabulaire
% moderne sous lequel chercher ce que les feuillets construisent. Le manifeste
% du site (`npm run manifest`) les extrait pour étiqueter la cote entière —
% cette ligne est la source unique des tags, il n'y a pas de fichier à côté.
\newcommand{\keywords}[1]{\par\smallskip\noindent
  {\footnotesize\textsc{Keywords} --- \textit{#1}}\par}

% --- L'appareil critique ----------------------------------------------------

% Le numéro de feuillet, en marge. C'est l'ancre qui permet de régler un
% désaccord sur une formule en regardant *un* feuillet plutôt qu'une chemise
% de six cents. Le site s'en sert aussi pour tourner les pages du fac-similé
% quand on fait défiler la transcription.
\newcommand{\leaf}[1]{%
  \marginnote{\footnotesize\color{gray!70}\textsf{#1}}%
  \label{leaf:#1}\ignorespaces}

% Illisible : on ne devine pas. Un mot inventé qui se lit comme les autres est
% le pire résultat possible.
\newcommand{\ill}{\textcolor{red!60!black}{[\,\ldots\,]}}

% Lecture douteuse : le mot est donné, et signalé comme incertain.
\newcommand{\uncertain}[1]{\uline{#1}}

% Ajout de l'éditeur — une lettre manquante, un mot sous-entendu.
\newcommand{\add}[1]{\textcolor{blue!50!black}{[#1]}}

% Biffé par Grothendieck lui-même. Ce qu'il a rayé fait partie du document :
% une rature dit souvent où la pensée a changé de direction.
\newcommand{\struck}[1]{\sout{#1}}

% Note du transcripteur, sur l'état matériel du feuillet.
\newcommand{\note}[1]{\par\smallskip{\small\color{gray!60!black}\textit{[#1]}}\par\smallskip}

% Note marginale de Grothendieck, distincte du corps du texte.
\newcommand{\marginal}[1]{\par\smallskip\noindent\hspace{1em}%
  {\small\color{gray!60!black}#1}\par\smallskip}

% --- Titre ------------------------------------------------------------------

\AtBeginDocument{%
  \begin{center}
    {\large\bfseries Fonds Alexandre Grothendieck --- cote n\textsuperscript{o}~\grothFolder}\\[2pt]
    {\itshape\grothFoldertitle}\\[4pt]
    {\small feuillets \grothFirst--\grothLast\ (lot \grothBatch)}\\[2pt]
    {\footnotesize\color{gray!60!black}Transcription établie d'après le fac-similé numérisé
     par l'Université de Montpellier.\\
     Les lectures incertaines sont soulignées, les illisibles notées [\,\ldots\,],
     les ajouts entre crochets.}
  \end{center}
  \vspace{1em}}

\endinput


\folder{8}
\batch{1}
\pages{1}{20}
\dating{[vers 1967- vers 1971]}
\watermark{Édition de démonstration}
\foldertitle{Cristaux (1970). Gribouillis et calculs divers : notes manuscrites (s.d.)}

\begin{document}

\page{1}

\marginal{encadré en haut à gauche : $C\colon \omega \to \omega^{(p)}$,
\quad \uncertain{$\Pi\colon \nu^{(p)} \to \nu$}}

\note{deux triangles dessinés à gauche, chacun « $\simeq$ » à un triangle
explicite à droite ; le sommet du second triangle de gauche est surchargé}

\begin{tikzcd}[column sep=small, row sep=small, nodes={font=\scriptsize}]
 & \check{\ell}^{\,G^*}_\bullet = \ell \arrow[dl, "(1)"'] & \\
 \ell^{G}_\bullet[-1] \arrow[rr] & & \Delta^*(G) \arrow[ul]
\end{tikzcd}

\qquad $\simeq$ \qquad

\begin{tikzcd}[column sep=small, row sep=small, nodes={font=\scriptsize}]
 & {[\omega \xrightarrow{C} \omega^{(p)}]} \arrow[dl, "{(\mathrm{id},0)}"'] & \\
 {[\omega^{(p)} \xrightarrow{0} \omega]} \arrow[rr, "{(\mathrm{id},C)}"'] & & {[\omega^{(p)} \xrightarrow{0} \omega^{(p)}]} \arrow[ul, "{(0,\mathrm{id})}"']
\end{tikzcd}

\note{la flèche $(\mathrm{id},0)$ porte « deg 1 » ; un mot biffé devant elle
est illisible}

\begin{tikzcd}[column sep=small, row sep=small, nodes={font=\scriptsize}]
 & \ell^{G(p)}_\bullet[1] \arrow[dl, "(1)"'] & \\
 \check{\ell}^{\,G^*(p)}_\bullet \arrow[rr] & & \Delta^*(G) \arrow[ul]
\end{tikzcd}

\qquad $\simeq$ \qquad

\begin{tikzcd}[column sep=small, row sep=small, nodes={font=\scriptsize}]
 & {[\omega^{(p^2)} \xrightarrow{0} \omega^{(p)}]} \arrow[dl, "{(\mathrm{id},0)}"'] & \\
 {[\omega^{(p)} \xrightarrow{C^{(p)}} \omega^{(p^2)}]} \arrow[rr, "{(\mathrm{id},0)}"'] & & {[\omega^{(p)} \xrightarrow{0} \omega^{(p)}]} \arrow[ul, "{(C,\mathrm{id})}"']
\end{tikzcd}

\note{le sommet du premier triangle de droite de la seconde ligne se lit
$\ell^{G(p)}_\bullet[1]$ sous une surcharge (« $= p^{(p)}$ ») ; sous la base
gauche, « $\| \; \ell^{(p)}$ ». La flèche $(\mathrm{id},0)$ porte de nouveau
« deg 1 »}

\[
\begin{array}{c}
S_{\mathrm{cris}} \\ \downarrow \varphi \\ S_{\mathrm{zar}}
\end{array}
\ \Big)\,\varepsilon
\qquad \varphi\varepsilon = F_S
\]
(frobenius absolu)

$\Delta^*(G)$ complexe cristallin avec $F$, $V$, filtration à facteurs
$\bigl(\varphi^*(Q), \varphi^*(P)\bigr)$ + filtration sur
$\Delta^*(G) = \varepsilon^*(\Delta(G))$ à facteurs $P$, $Q$.

\[
\varphi^*(\omega) \xrightarrow{\ \varphi^*(C)\ } \varphi^*(\omega^{p}),
\qquad \varphi^*(\omega) = \mathbb{D}^*(G), \quad
\varphi^*(\omega^{p}) = \mathbb{D}^*(G)^{(p)},
\qquad V \downarrow, \quad F = 0 \leftarrow
\]
\note{schéma en carré : $V$ descend de $\varphi^*(C)$, et une flèche
« $F = 0$ » revient de $\mathbb{D}^*(G)^{(p)}$ vers $\mathbb{D}^*(G)$}

\[
\mathcal{M} \underset{F}{\overset{V}{\rightleftarrows}} \mathcal{M}^{(p)},
\qquad \mathcal{M} = \mathbb{D}^*(G),
\]
\struck{$\Delta^*(G) = \varepsilon^*(\mathcal{M})^{\circ}$}
\[
\Delta^*(G) = [\varepsilon^*(\mathcal{M}) \xrightarrow{0} \varepsilon^*(\mathcal{M})]
= \varepsilon^*[\mathcal{M} \to \mathcal{M}],
\]
avec une flèche de $\mathcal{M} \rightleftarrows \mathcal{M}^{(p)}$ vers
$[\varepsilon^*(\mathcal{M}) \xrightarrow{0} \varepsilon^*(\mathcal{M})]$.

\begin{tikzcd}[column sep=small, row sep=small, nodes={font=\scriptsize}]
 & {[\mathcal{M}^{(p)} \xrightarrow{V} \mathcal{M}^{(p)}]} \arrow[dl, "{(\mathrm{id},0)}"'] & \\
 {[\mathcal{M}^{(p)} \xrightarrow{F} \mathcal{M}^{(p)}]} \arrow[rr, "{(\mathrm{id},0)}"'] & & {\Delta^*(G) = [\mathcal{M} \xrightarrow{0} \mathcal{M}]} \arrow[ul, "{(V,\mathrm{id})}"']
\end{tikzcd}

\note{les exposants des deux crochets sont surchargés ; la lecture
$\mathcal{M}^{(p)}$ aux deux places est douteuse. La flèche $(\mathrm{id},0)$
de gauche porte « deg 1 »}

\[
0 \to {}_F\mathcal{M}^{(p)} \to \mathcal{M} \to {}_V\mathcal{M}^{(p)} \to
\]

\page{3}

\struck{\ill{}} \uncertain{gpes} finis loc.\ libres annulés par $p$ en car.\ $p$.

a) Cristal (loc.\ libre de \uncertain{type fini ?}) $\mathcal{M}$ sur
$S_{\mathrm{cris}}/\mathbb{F}_p$ avec
\[
\mathcal{M}^{(p)} \underset{V_{\mathcal{M}}}{\overset{F_{\mathcal{M}}}{\rightleftarrows}} \mathcal{M}
\]
satisfaisant $FV = VF = 0$.

On en déduit complexe cristallin parfait, d'amplitude parfaite
$\subset [0,1]$,
\[
\Delta^* = \Delta^*(\mathcal{M}), \qquad
\mathcal{H}^0(\Delta^*) \simeq \mathcal{H}^1(\Delta^*) \simeq \mathcal{M}
\]
avec $F_{\Delta^*}$, $V_{\Delta^*}$ satisfaisant
$F_{\Delta^*}V_{\Delta^*} = 0$, $V_{\Delta^*}F_{\Delta^*} = 0$, et triangles
exacts

\begin{tikzcd}[column sep=small, row sep=small]
 & {[\mathcal{M} \xrightarrow{V} \mathcal{M}^{(p)}]} \arrow[dl, "{\partial_{\mathrm{I}},\ 1}"'] & \\
 {[\mathcal{M}^{(p)} \xrightarrow{F} \mathcal{M}]} \arrow[rr, "i_{\mathrm{I}}"'] & & \Delta^{*(p)} \arrow[ul, "\pi_{\mathrm{I}}"']
\end{tikzcd}

\quad $(\mathrm{I})_{\mathrm{cris}}$, \qquad
$\Delta^{*(p)} \simeq \varphi^*(\Delta^*)$ \quad
$(\Delta^* = \varepsilon^*(\Delta^*))$

\begin{tikzcd}[column sep=small, row sep=small]
 & {[\mathcal{M}^{(p)} \xrightarrow{F} \mathcal{M}]} \arrow[dl, "{\partial_{\mathrm{II}},\ 1}"'] & \\
 {[\mathcal{M} \xrightarrow{V} \mathcal{M}^{(p)}]} \arrow[rr, "i_{\mathrm{II}}"'] & & \Delta^* \arrow[ul, "\pi_{\mathrm{II}}"']
\end{tikzcd}

\quad $(\mathrm{II})_{\mathrm{cris}}$

\note{la flèche horizontale de $(\mathrm{I})$ est notée $\partial_{\mathrm{I}}$
ou $i_{\mathrm{I}}$ selon l'endroit ; on garde $i_{\mathrm{I}}$}

\marginal{N.B. Les complexes
$\ell^{\,\cdot}_\bullet(\mathcal{M}) = [\mathcal{M}^{(p)} \xrightarrow{F} \mathcal{M}]$ ($F$ nul),
$\ell^{\,\prime}_\bullet(\mathcal{M}) = (\mathcal{M} \xrightarrow{V} \mathcal{M}^{(p)})$ ($V$ nul),
$\Delta^*(\mathcal{M})$,
\ill{} avec $(F, V)$ \ill{}, les triangles exacts $(\mathrm{I})$, $(\mathrm{II})$
sont \uncertain{compatibles} avec $(F,V)$-structure}
\note{encadré à droite, dont le début (« N.B. ») et les lignes suivantes sont
serrés ; « ($F$ nul) », « ($V$ nul) » écrits en marge des deux premières
lignes}

b) Complexes parfaits $\ell_\bullet$, $\ell'_\bullet$ sur $S_{\mathrm{zar}}$,
\ill{} d'amplitude parfaits $\subset [-1,0]$,
\struck{triangle exact \ill{}} (III) :

\begin{tikzcd}[column sep=small, row sep=small]
 & \check{\ell}'_\bullet \arrow[dl, "\partial_0"'] & \\
 \ell_\bullet[-1] \arrow[rr, "i_0"'] & & \Delta^* \arrow[ul, "\pi_0"']
\end{tikzcd}

\quad $\Delta^* \overset{\mathrm{df}}{=} \varepsilon^*(\Delta^*)$

\note{au-dessus de b), entouré : « \ill{} $(F,V)$-\ill{} », et « d'amplitude »
repris à côté}

\[
\varphi^*(\ell_\bullet[-1]) \overset{\alpha}{\simeq} [\mathcal{M}^{(p)} \xrightarrow{F} \mathcal{M}]
\qquad
\varphi^*(\check{\ell}'_\bullet) \simeq [\mathcal{M} \xrightarrow{V} \mathcal{M}^{(p)}]
\]
fournissant un \struck{\ill{}} isom.\ de $(\mathrm{I}_{\mathrm{cris}})$ avec
$\varphi^*(\mathrm{III})$

\note{le passage suivant est entre crochets}
[d'où, posant $M = \varepsilon^*(\mathcal{M})$ loc.\ libre sur
$S_{\mathrm{zar}}$, avec connexion à courbure nulle,
\[
\ell_\bullet[-1]^{(p)} \simeq [M^{(p)} \xrightarrow{F_M} M],
\qquad
\check{\ell}'_\bullet{}^{(p)} \simeq [M \xrightarrow{V_M} M^{(p)}] \quad ! \;]
\]
\marginal{comp.\ avec connexion}

\emph{Compatibilités} : \marginal{Posons}
\[
\begin{cases}
\pi_0\, \varepsilon^*(i_{\mathrm{II}}) \text{ ``$=$'' } F_{\check\ell'_\bullet} = {}^tV_{\ell'_\bullet} \\
\varepsilon^*(\pi_{\mathrm{II}})\, i_0 \text{ ``$=$'' } V_{\ell_\bullet}[-1]
\end{cases}
\]

\page{4}

\note{feuillet écrit tête-bêche par rapport au scan ; on le lit retourné. Deux
longs traits de crayon le traversent en biais sans le biffer}

$G$ \uncertain{gpe} fini loc.\ libre, annulé par $p$, en car.\ $p>0$ ;
on lui associe (conjecturalement) :

a) $\mathbb{D}^*(G)$ Module (de présentation finie ?
\uncertain{loc.\ libre de type fini ?}) sur $S_{\mathrm{cris}}$ (topos
cristallin de Berthelot de $S$ sur $\mathbb{F}_p$)
[\uncertain{fonctoriel} en $G$, \emph{\uncertain{donc}}]
\[
\mathbb{D}^*(G)^{(p)} \underset{V}{\overset{F}{\rightleftarrows}} \mathbb{D}^*(G)
\]
avec $FV = 0$, $VF = 0$

b) homomorphismes
\note{« b) homomorphismes » reste sans suite ; la ligne suivante, serrée
contre elle, commence par c)}
c) des complexes $\ell^{G}_\bullet$ et $\ell^{G^*}_\bullet$ sur $S$, parfaits
d'amplitude parfaits dans $[-1,0]$, et un hom
\[
(\ell^{G^*}_\bullet)^{\vee} \xrightarrow{\ \pi\ } \ell^{G}_\bullet
\]
i.e.\ un hom
\[
\underline{O}_S \longrightarrow \ell^{G^*}_\bullet \overset{L}{\otimes} \ell^{G}_\bullet
\]

d) \struck{Deux} Un triangle\uncertain{s} exact\uncertain{s} sur $S_{\mathrm{cris}}$

\begin{tikzcd}[column sep=small, row sep=small]
 & \varphi^*(\ell^{G^*}_\bullet)^*[-1] \arrow[dl, "\partial"'] & \\
 \varphi^*(\check{\ell}^{\,G^*}_\bullet) \arrow[rr] & & \Delta^*_{\mathrm{cris}}(G) \arrow[ul]
\end{tikzcd}

\qquad
$\bigl[\,\varphi^*(\pi)\colon \underline{O}_{S_{\mathrm{cris}}} \xrightarrow{\ 2\ }
\varphi^*(\check{\ell}^{\,G}_\bullet \otimes \check{\ell}^{\,G^*}_\bullet)\,\bigr]$

(qui détermine donc $\Delta^*_{\mathrm{cris}}(G)$ : isom.\ (à \uncertain{un
unique} près) et des isom.
\[
\mathcal{H}^0(\Delta^*_{\mathrm{cris}}(G)) \simeq \mathbb{D}^*(G), \qquad
\mathcal{H}^1(\Delta^*_{\mathrm{cris}}(G)) \simeq \mathbb{D}^*(G) .
\]

(d'où par application de $\varepsilon^*$ un triangle

\begin{tikzcd}[column sep=small, row sep=small]
 & (\ell^{G}_\bullet)^{(p)}[-1] \arrow[dl, "\varepsilon^*(\partial)"'] & \\
 (\check{\ell}^{\,G^*}_\bullet)^{(p)} \arrow[rr] & & \Delta^*(G) \arrow[ul]
\end{tikzcd}

\qquad $\Delta^*(G) = \varepsilon^*(\Delta^*_{\mathrm{cris}}(G))$

e) Un triangle exact
\note{la page s'arrête là}

\page{5}

\[
\begin{array}{ll}
V_{\ell_\bullet}\colon \ell_\bullet \to \ell^{(p)}_\bullet &
F_{\ell_\bullet} = 0\colon \ell^{(p)}_\bullet \to \ell_\bullet \\
V_{\ell'_\bullet}\colon \ell'_\bullet \to \ell'^{(p)}_\bullet &
F_{\ell'_\bullet} = 0\colon \ell'^{(p)}_\bullet \to \ell'_\bullet
\end{array}
\]

1) Le triangle (III) est compatible avec $(F,V)$-structures : \emph{à
prouver} !

2) L'isom.\ $\varphi^*(\mathrm{III}) = \mathrm{I}$ compatible avec
$(F,V)$-structures.
\note{« isom. » est souligné d'un trait ondulé}

1) \struck{\ill{}}

\begin{tikzcd}[column sep=small, row sep=small, nodes={font=\scriptsize}]
\Delta^{*p} \arrow[rr, "F_{\Delta^*}"] \arrow[dr, "\pi_0^{(p)}"'] & & \Delta^* \\
 & \check\ell'^{(p)}_\bullet \arrow[ur, "\varepsilon^*(i_{\mathrm{II}})"'] &
\end{tikzcd}

\note{sous $\pi_0^{(p)}$ :
« $\varepsilon^*(\varphi^*(\pi_\cdot)) = \varepsilon^*(\pi_{\mathrm{I}})$ »}
\marginal{implique que $i_0$ comp.\ avec $F$}

\begin{tikzcd}[column sep=small, row sep=small, nodes={font=\scriptsize}]
\Delta^* \arrow[rr, "V_{\Delta^*}"] \arrow[dr, "\varepsilon^*(\pi_{\mathrm{II}})"'] & & \Delta^{*(p)} \\
 & \ell^{(p)}_\bullet[-1] \arrow[ur, "{i_0^{(p)}}"'] &
\end{tikzcd}

\qquad
$i_0^{(p)} = \varepsilon^*\varphi^*(i_0) = \varepsilon^*(i_{\mathrm{I}})$
\marginal{implique que $\pi_0$ compatible avec $V$}

$i_0$ compatible avec $V$ :

\begin{tikzcd}[column sep=small, row sep=small, nodes={font=\scriptsize}]
\ell_\bullet[-1] \arrow[r, "i_0"] \arrow[d, "{V_{\ell_\bullet}[-1]}"'] & \Delta^* \arrow[d, "V_{\Delta^*}"] \arrow[dr, "\varepsilon^*(\pi_{\mathrm{II}})"] & \\
\ell_\bullet[-1]^{(p)} \arrow[r, "i_0^{(p)}"'] & \Delta^{*(p)} & \ell^{(p)}_\bullet[-1] \arrow[l, "i_0^{(p)}"]
\end{tikzcd}

\quad OK.

\[
\varepsilon^*(\pi_{\mathrm{II}})\, i_0 \ne V_{\ell_\bullet}[-1]
\]
\note{le signe entre les deux membres est une égalité barrée ; la lecture
« $\ne$ » reste douteuse}

On voit de même que $\pi_0$ compatible avec $F$.

$\partial_0$ compatible avec $F$ :

\begin{tikzcd}[column sep=small, row sep=small, nodes={font=\scriptsize}]
\check\ell'^{(p)}_\bullet \arrow[r, "{\pi_0\,\varepsilon^*(i_{\mathrm{II}})}"] \arrow[d, "\partial_0^{(p)}"'] & \check\ell'_\bullet \arrow[d, "\partial_0"] \\
\ell_\bullet[-1]^{(p)} \arrow[r, "0"'] & \ell_\bullet[-1]
\end{tikzcd}

\qquad
$\partial_0\, \pi_0\, \varepsilon^*(i_{\mathrm{II}}) = 0$

On voit de même $\partial_0$ compatible avec $V$.

2) À prouver que
\[
\begin{cases}
\varphi^*(V_{\ell_\bullet}[-1]) = V_{[\mathcal{M}^{(p)} \xrightarrow{F} \mathcal{M}]} \\
\varphi^*(F_{\check\ell'_\bullet}[-1]) = F_{[\mathcal{M} \xrightarrow{V} \mathcal{M}^{(p)}]}
\end{cases}
\]
\[
\varphi^*(V_{\ell_\bullet}[-1]) = \varphi^*(\varepsilon^*(\pi_{\mathrm{II}})\, i_0)
= \varphi^*\varepsilon^*(\pi_{\mathrm{II}})\, \varphi^*(i_0)
= \bigl[\pi_{\mathrm{II}}^{(p)}\, i_{\mathrm{I}}
= V_{[\mathcal{M}^{(p)} \to \mathcal{M}]}\bigr]
\]
et l'autre relation se \uncertain{prouve} de même !!

\page{6}

\note{feuillet de tapuscrit dactylographié, numéroté 35 à l'encre en haut à
droite, sans rapport avec les notes manuscrites voisines ; il est gardé parce
qu'il porte des reprises à la main, et le texte tapé n'est pas recomposé. Il
achève une démonstration : puisque (5.1)(iii) est vrai pour
$(g, \bar f, \bar G)$ et pour $(g', \pi, \bar G)$, les faisceaux
$R^q\bar f'_*(\bar g^*\bar G)$ et $R^q\pi'_*(\bar g^*\bar G)$ sont nuls pour
$q>0$ ; la suite spectrale de Leray
$E_2^{p,q} = R^pf'_*(R^q\pi'_*(\bar g^*\bar G)) \Rightarrow R^n\bar f'_*(\bar G)$
donne $R^pf'_*(\pi'_*\bar g^*\bar G) = 0$ pour $p>0$, et l'injection
$\pi_*\bar F \to \pi_*\bar G$ est un effacement des foncteurs
$(R^pf'_*)g'^*$ dans la catégorie des faisceaux de torsion. Puis « Lemme
(8.4) : Il suffit de vérifier le théorème (5.1)(i) (resp. (i) et (ii), resp.
(iii)) pour $(g, f, F)$ avec $f$ projectif », et le début de sa démonstration
(réduction, d'après (6.5) avec $P$ la propriété vide, au cas $S$ strictement
local et $g$ l'inclusion du point fermé ; famille $\mathfrak{F}$ des
sous-préschémas fermés $h\colon Y \to X$ contenant la fibre fermée $X_0$ ;
morphisme $H^q(X,F) \to H^q(X, h_*h^*F)$). Les parenthèses, les étoiles et
plusieurs lettres des formules sont repassées à l'encre. Interventions à la
main : « s'ensuit » biffé et remplacé au-dessus par « résulte » ;
après « torsion » un renvoi entre parenthèses, lu « (IX 1.2 (v)) » (lecture
douteuse) ; un crochet double en marge droite à la fin du paragraphe ;
« applicable » biffé et remplacé au-dessus par « pertinent, »}

\page{7}

\begin{tikzcd}[column sep=small, row sep=small, nodes={font=\scriptsize}]
\mathcal{M} \arrow[r, "V"] \arrow[d, "V"'] & \mathcal{M}^{(p)} \arrow[d, "V^{(p)}"] \\
\mathcal{M}^{(p)} \arrow[r, "V^{(p)}"'] & \mathcal{M}^{(p^2)}
\end{tikzcd}

\note{une flèche $F$ remonte à gauche, et
${}^tF$ \uncertain{} est inscrit au milieu ; à côté : « $V^{(p)}F$ »}

$\mathcal{M}$

\begin{tikzcd}[column sep=small, row sep=small]
 & B \arrow[dl, "\partial_{\mathrm{I}}"'] & \\
 A \arrow[rr, "i_{\mathrm{I}}"'] & & \mathfrak{X}^p \arrow[ul, "\pi_{\mathrm{I}}"']
\end{tikzcd}

\quad (I) \qquad

\begin{tikzcd}[column sep=small, row sep=small]
 & A \arrow[dl, "\partial_{\mathrm{II}}"'] & \\
 B \arrow[rr, "i_{\mathrm{II}}"'] & & \mathfrak{X} \arrow[ul, "\pi_{\mathrm{II}}"']
\end{tikzcd}

\quad (II)

\begin{tikzcd}[column sep=small, row sep=small]
 & A^p \arrow[dl, "\partial_{\mathrm{II}}^{(p)}"'] & \\
 B^p \arrow[rr, "i_{\mathrm{II}}^{(p)}"'] & & \mathfrak{X}^p \arrow[ul, "\pi_{\mathrm{II}}^{(p)}"']
\end{tikzcd}

\quad $(\mathrm{II}^{(p)})$

\note{encadré à gauche :}
\[
\begin{array}{ll}
\pi_{\mathrm{I}} : & \mathfrak{X}^{p} \underset{V_{\mathfrak{X}}}{\overset{F_{\mathfrak{X}}}{\rightleftarrows}} \mathfrak{X} \\
i_{\mathrm{I}}\,\pi_{\mathrm{II}} : & \mathfrak{X} \xrightarrow{V_{\mathfrak{X}}} \mathfrak{X}^{p}
\end{array}
\]
\note{le premier membre de gauche, surchargé, se lit
$\pi_{\mathrm{II}}^{(p)}\, i_{\mathrm{I}}$ ou $\pi_{\mathrm{I}}$}

\[
\left\lbrace
\begin{array}{lll}
\partial_{\mathrm{I}}\,\partial_{\mathrm{II}} : & A^{2} \to A & (\mathrm{nul}\,?) \\
\partial_{\mathrm{II}}\,\partial_{\mathrm{I}} : & B^{2} \to B & (\mathrm{nul}\,?) \\
\pi_{\mathrm{II}}^{(p)}\,\partial_{\mathrm{I}} : & A \xrightarrow{V_A} A^{(p)} & \\
\pi_{\mathrm{I}}\,, i_{\mathrm{II}}^{(p)} : & B^{p} \xrightarrow{F_B} B &
\end{array}
\right.
\]
\note{« nul ? » entouré deux fois ; les exposants « 2 » sont écrits sur des
$p$ ; une accolade relie ces lignes à « $\mathcal{M}$ » plus bas}

\begin{tikzcd}[column sep=small, row sep=small, nodes={font=\scriptsize}]
 & \varphi^*(b) \arrow[dl, "{\partial_{\mathrm{I}} = \varphi^*(\partial_0)}"'] & \\
 \varphi^*(a) \arrow[rr, "{i_{\mathrm{I}} = \varphi^*(i_0)}"'] & & \mathfrak{X}^p \arrow[ul, "{\pi_{\mathrm{I}} = \varphi^*(\pi_0)}"']
\end{tikzcd}

\quad $\varphi^*(\mathrm{III}) \simeq (\mathrm{I})$
\qquad

\begin{tikzcd}[column sep=small, row sep=small]
 & b \arrow[dl, "\partial_0"'] & \\
 a \arrow[rr, "i_0"'] & & \mathfrak{X} \arrow[ul, "\pi_0"']
\end{tikzcd}

\quad III

\[
\begin{array}{ccccccc}
0 & \to & \mathcal{M} & \to & \mathcal{M}^p & \to & \cdot \\
 & & \downarrow & & \downarrow & & \\
 & \mathcal{M}^p & \to & \mathcal{M} & \to & 0 & \to \\
 & & \downarrow & & \downarrow & & \\
 & \mathcal{M} & \to & \mathcal{M}^p & \to & 0 & \to
\end{array}
\]
\note{esquisse de gauche, rattachée par une accolade à l'encadré ; les zéros
et les flèches sont incertains}

$\varepsilon^*(\mathrm{I}) \simeq \mathrm{II}^{(p)}$ :

\begin{tikzcd}[column sep=small, row sep=small]
 & b^p \arrow[dl, "\partial_0^{(p)}"'] & \\
 a^p \arrow[rr, "i_0^{(p)}"'] & & X^p \arrow[ul, "\pi_0^{(p)}"']
\end{tikzcd}

$\varepsilon^*(\mathrm{II})$ :

\begin{tikzcd}[column sep=small, row sep=small]
 & a^p \arrow[dl, "\varepsilon^*(\partial_{\mathrm{II}})"'] & \\
 b^p \arrow[rr, "\varepsilon^*(i_{\mathrm{II}})"'] & & X \arrow[ul, "\varepsilon^*(\pi_{\mathrm{II}})"']
\end{tikzcd}

\marginal{triangles compatibles avec $F$-$V$ structures}

\[
\begin{array}{ll}
\mathfrak{X}^p \to B \to \mathfrak{X} &
\mathfrak{X} \to A \to \mathfrak{X}^{(p)} \\
X^{(p)} \to b^{(p)} \to X &
X \to a^{(p)} \to X^{(p)}
\end{array}
\]

\note{encadré du bas :}
\[
\left\lbrace
\begin{array}{ll}
\pi_0\, \varepsilon^*(i_{\mathrm{II}}) : & b^{(p)} \xrightarrow{F_b} b \quad ? \\
\varepsilon^*(\pi_{\mathrm{II}})\, i_0 : & a \xrightarrow{V_a} a^{(p)} \quad ?
\end{array}
\right.
\]

\page{8}

\note{feuillet du même tapuscrit, numéroté 37.1 à l'encre. Suite d'une
démonstration : les flèches verticales (d'un diagramme sur la page
précédente du tapuscrit) sont injectives (6.6) ; $\alpha_0 \in H^0(X_0,F_0)$
se relève à $H^0(X,F)$ si et seulement si son image se relève dans
$H^0(X_0,G_0)$, avec $G = (h_*h^*F) \times (\pi_*\pi^*F)$ ; appliquant
(5.1)(i) à $(g,\pi,\cdot)$ et à $(g,\bar f,\cdot)$ on conclut que
$H^0(X, \pi_*\pi^*F) \xrightarrow{\sim} H^0(X_0, (\pi_*\pi^*F)_0)$, ce qui
achève la démonstration dans ce cas. Puis, supposant (5.1)(i) vrai et
(5.1)(ii) vrai pour chaque morphisme projectif, pour $F$ un faisceau de
groupes ind-finis sur $X$, les applications
$H^1(X,F) \xrightarrow{u} H^1(Y,h^*F) \xrightarrow{v} H^1(X_0,F_0)$, avec
$vu$ injectif (6.1), et la surjectivité de $u$ à démontrer ; « d'après
(8.2(ii)) appliqué à $f_Y\colon Y \to S$ », injection $h^*F \to G^*$ ;
« Soit $j\colon X' \to X$ l'image (réduite) de $\pi\colon \bar X \to X$ ».
Interventions à la main : en marge gauche, « \uncertain{un peu bref,} faire
l'argument » (lecture douteuse), en regard de « on conclut que », souligné ;
à la fin du paragraphe, « (5.1.(i)) » ajouté ; « (8.2 (ii)) » écrit à
l'encre sur un renvoi tapé ; entourée à la suite de « soit nulle », la
parenthèse « (avec la structure induite) », rattachée par un trait à
« (réduite) » ; « des » ajouté au-dessus de « fermé » ; nombreuses lettres
et accents repassés}

\page{9}

\uncertain{Groupes} loc.\ lib.\ en car.\ $p$ (\uncertain{par ex.\ annulés par $p$})

a) Cristal de prés.\ finie sur $S_{\mathrm{cris}}/\struck{\ill{}}\,\mathbb{Z}_p$
\struck{(par ex.\ \ill{} De \ill{})}
d'amplitude parfaite $\subset [-1,0]$, soit $\mathcal{M}$, avec structures
$F_{\mathcal{M}}$ et $V_{\mathcal{M}}$.

b) \struck{\ill{}}

On en déduit deux triangles exacts de complexes avec $F$, $V$ sur
$S_{\mathrm{cris}}/\mathbb{F}_p$

\begin{tikzcd}[column sep=small, row sep=small, nodes={font=\scriptsize}]
 & \mathcal{M}\overset{L}{\otimes}/V \arrow[dl, "\partial_{\mathrm{I}}"'] & \\
 \mathcal{M}\overset{L}{\otimes}/F \arrow[rr, "i_{\mathrm{I}}"'] & & \mathcal{M}^{(p)}\overset{L}{\otimes}/p \arrow[ul, "\pi_{\mathrm{I}}"']
\end{tikzcd}

\quad (I), \qquad
$\mathcal{M}^{(p)}\overset{L}{\otimes}/p \simeq \varphi^*(\Delta^*)$ \quad
$\Delta^* = \varepsilon^*(\Delta^*)$

\begin{tikzcd}[column sep=small, row sep=small, nodes={font=\scriptsize}]
 & \mathcal{M}\overset{L}{\otimes}/F \arrow[dl] & \\
 \mathcal{M}\overset{L}{\otimes}/V \arrow[rr] & & \mathcal{M}\overset{L}{\otimes}/p = \Delta^* \arrow[ul]
\end{tikzcd}

\quad (II)

\note{sous le sommet droit de (I), une expression biffée « \ill{}
$(\mathcal{M})$ » ; sous (II) : « $\simeq^{\mathrm{dbs}}$
$\overset{L}{\otimes}\,\Delta^*(\mathcal{M})$ » (lecture douteuse)}

$\alpha\colon S_{\mathrm{cris}}/\mathbb{F}_p \hookrightarrow S_{\mathrm{cris}}/\mathbb{Z}_p$

b) \struck{\ill{}} Descente à $S_{\mathrm{zar}}$ (via $\varphi$) de la
« filtration » (I) sur $\varphi^*(\Delta^*)$, ou un triangle exact

\begin{tikzcd}[column sep=small, row sep=small]
 & \check\ell'_\bullet \arrow[dl, "\partial_0"'] & \\
 \ell_\bullet[-1] \arrow[rr, "i_0"'] & & \Delta^* \arrow[ul, "\pi_0"']
\end{tikzcd}

Prolongement à un \uncertain{épaississement} à puiss.\ \uncertain{div.}\ $S'$ de
$S$ (\uncertain{i.e.} avec $\mathfrak{p} = \underline{\mathrm{Var}}(p.\,\mathcal{I}_S)$)
On considère
\[
\mathcal{M} \overset{\mathbb{L}}{\otimes} \mathcal{O}_{S'} = \Delta^*(S'),
\]
qui prolonge $\Delta^*$, et il faut se donner un prolongement de la filtration
b) de $\Delta^*$ en

\begin{tikzcd}[column sep=small, row sep=small]
 & \bar{\check\ell}'_\bullet \arrow[dl, "\bar\partial_0"'] & \\
 \bar\ell_\bullet[-1] \arrow[rr, "\bar i_0"'] & & \Delta^*(S') \arrow[ul, "\pi_0"']
\end{tikzcd}

\page{10}

\note{feuillet du même tapuscrit, écrit tête-bêche par rapport au scan,
numéroté 32 à l'encre ; deux flèches verticales au crayon y sont tracées, sans
rapport lisible avec les notes. Texte : l'assertion (iii) se traite de façon
analogue ; les foncteurs $g^*(R^qf_*)$ étant effaçables pour $q>0$ dans la
catégorie des faisceaux de torsion, la bijectivité de $\varphi^n$ pour $F$
implique l'effaçabilité de $(R^nf'_*)g'^*$ pour $F$ ; inversement, supposant
$\varphi^{n-1}$ bijectif pour $(g,f,\cdot)$ et $(R^nf'_*)g'^*$ effaçable pour
$F$, et $F \to G$ un effacement de $(R^{n-1}f'_*)g'^*$ avec $G$ de torsion,
$C = G/F$, on complète le diagramme $(*)$ de deux suites exactes longues
reliées par $\varphi^{n-1}$, $\varphi^{n}$. Interventions à la main : le
renvoi « (8.1) » écrit à l'encre dans une parenthèse laissée vide ; les
$\varphi$ encrés, cédilles et accents ajoutés}

\page{11}

a) Les catégories de structures a) b) introduites précédemment
définissent-elles des \emph{champs} sur $S_{\mathrm{zar}}$ ?

b) Le foncteur « \uncertain{restriction} partiel » (de $\mathcal{C}$ dans b))
allant vers les cristaux $\mathcal{M}$ sur $S_{\mathrm{cris}}/\mathbb{F}_p$
est-il \emph{fidèle} ?

c) Le foncteur naturel \uncertain{Mais} des modules loc.\ libres $\omega$ sur
$S_{\mathrm{zar}}$, avec $F_\omega$, $V_\omega$, vers les structures a), b),
est-il pleinement fidèle (NB il est fidèle) ; l'est-il du moins
\struck{si $F_\omega$} sur la sous-catégorie définie par $F_\omega = 0$ ;
l'est-il sur la \uncertain{sous-catégorie} où $\operatorname{Ker} F = \operatorname{Im} V$ ?

d) Peut-on à $(\omega, F_\omega, V_\omega)$ associer un schéma en groupes
fini loc.\ libre sur $S$ ?

\marginal{\ill{} que les \struck{\ill{}}
\struck{$H^0(\mathcal{C}) \to H_0(\check\ell'[1]) \to \cdots$}
\ill{} $\cdots \to H_1(\ell) \to H_1(\check\ell)$ \struck{\ill{}}
$0 \to t^* \to \omega$ nul, et $\omega^*$ et $t^*$ loc.\ libres}
\note{note écrite en biais dans la marge gauche, en regard de b) ; en partie
biffée et encadrée, lue par fragments}

\page{12}

Cas des groupes de BT tronqués annulés par $p$.

Signifie que $F_{\mathcal{M}}$, $V_{\mathcal{M}}$ forment suite exacte
$\Longleftrightarrow$ $F_M$, $V_M$ forment suite exacte \ldots

Alors la donnée b) implique, passant aux $\underline{H}^1$
\[
\Bigl[\omega = \underline{H}^1(\ell_\bullet[-1]) = \underline{H}^0(\ell_\bullet),
\quad \nu' = \underline{H}^1(\check\ell'_\bullet)\Bigr]
\]
($\omega$, $\nu'$ loc.\ libres)
$M$ ($= \varepsilon^*(\mathcal{M})$) filtré par
\[
0 \to \omega \to M \to \nu' \to 0
\]
de façon que dans
\[
0 \to \varphi^*(\omega) \to \varphi^*(M) \to \varphi^*(\nu') \to 0, \qquad
\varphi^*(M) = \mathcal{M}^{(p)},
\]
$\varphi^*(\omega)$ ne soit autre que
$\operatorname{Im}(\mathcal{M} \xrightarrow{V_{\mathcal{M}}} \mathcal{M}^{(p)})$.

\marginal{Cela détermine $\omega$ de façon \emph{unique} !}

De même, passant aux $\underline{H}^0$, on trouve
\[
\Bigl[n = \underline{H}^0(\ell_\bullet[-1]) = \underline{H}^{-1}(\ell_\bullet),
\quad t' = \underline{H}^0(\check\ell'_\bullet)\Bigr]
\]
\[
0 \to n \to M \to t' \to 0
\]
de façon que dans
\[
0 \to \varphi^*(n) \to \varphi^*(M) \to \varphi^*(t') \to 0
\]
on ait
\[
\varphi^*(n) = \operatorname{Ker}(\mathcal{M}^{(p)} \xrightarrow{F_{\mathcal{M}}} \mathcal{M}) .
\]
Comme $\operatorname{Ker} F_{\mathcal{M}} = \operatorname{Im} V_{\mathcal{M}}$,
on conclut
\[
\varphi^*(n) = \varphi^*(\omega),
\]
d'où
\[
\boxed{n = \omega}, \quad \boxed{t' = \nu'} .
\]

La seule donnée de \struck{\ill{}} $\mathcal{M}$, $F_{\mathcal{M}}$,
$V_{\mathcal{M}}$ ($\Rightarrow \omega \subset M = \varphi^*(\mathcal{M})$)
\uncertain{permet-elle de} reconstituer toute la donnée b) ? On voit qu'on
pourra construire
\[
\omega \xrightarrow{V_\omega} \omega^{(p)}, \qquad
t'^{(p)} \xrightarrow{F_{t'}} t' \qquad
\bigl(\omega' = \check t', \ \omega' \xrightarrow{V_{\omega'}} \omega'^{(p)},
\ V_{\omega'} = {}^tF_{t'}\bigr)
\]
Car :
\[
\omega^{(p)} = \operatorname{Im}(M \xrightarrow{V_M} M^{(p)}),
\]
donc $V_M/\omega$ \ldots
\note{la fin de la ligne se lit « $V_M|\omega$ \ill{} $\omega$ dans
$\omega^{(p)}$ », douteux}
\[
t'^{(p)} = \operatorname{Coker}(M \xrightarrow{V_M} M^{(p)}),
\]
donc $F_M\colon M^{(p)} \to M$ se factorise par $t'^{(p)} \hookrightarrow M$,
d'où $F_{t'}\colon t'^{(p)} \hookrightarrow M \to t'$.

\note{en bas à gauche, un petit diagramme griffonné et biffé :
$M \twoheadrightarrow \omega$, $V_M$, $F_\omega$, $M^{(p)}$,
$M \twoheadrightarrow t$, $M^{(p)} \twoheadrightarrow t$ ; illisible dans le
détail}

\page{13}

\begin{tikzcd}[column sep=small, row sep=small]
M \arrow[r, two heads] \arrow[d, "V_M"'] \arrow[dr, "\mathrm{\acute{e}pi}", two heads] & \omega \arrow[d, "V_\omega"] \\
M^p \arrow[r, two heads] & \omega^{(p)}
\end{tikzcd}

\qquad

\begin{tikzcd}[column sep=small, row sep=small]
M^p \arrow[r, two heads] & t' \\
M^{(p)} \arrow[u] \arrow[r] & t'^{(p)} \arrow[u, "F_t"'] \arrow[ul, hook]
\end{tikzcd}

\struck{On peut dans \uncertain{former}} On aura

\begin{tikzcd}[column sep=small, row sep=small]
 & t^p \arrow[d] \arrow[dr, "F_{t'}"] & \\
0 \arrow[r] & M \arrow[r] \arrow[d] & t' \arrow[r] & 0 \\
\omega \arrow[ur] \arrow[r, "V_\omega"'] & \omega^p \arrow[d] & & \\
 & 0 & &
\end{tikzcd}

\note{l'esquisse est incomplète : au-dessus de $t^p$ un pointillé vertical ;
à gauche de $t^p$ un mot biffé et noirci ; la ligne horizontale se lit
« $0 - \omega - M - t' - 0$ », sans pointes de flèches, et la place de
$\omega$ dans le tikzcd n'est qu'approchée}

\page{14}

\note{feuillet du même tapuscrit, numéroté 33 à l'encre. Fin de la
démonstration précédente (injectivité et surjectivité de $\varphi^n$ par un
effacement $F \to G$ de $(R^nf'_*)g'^*$ avec $G$ de torsion, $C = G/F$, et le
diagramme $(*)$ complété par un $0$) ; puis « Lemme (8.3) : Soit
$X \xleftarrow{\pi} \bar X$, $f\colon X \to S$, $\bar f\colon \bar X \to S$
un diagramme commutatif de morphismes propres. Supposons que $\pi$ soit
surjectif », changement de base $g\colon S' \to S$, $g'\colon X' \to X$,
$\bar g\colon \bar X' \to \bar X$ ; (i), (ii), (iii) comme dans (5.1) ; le
début de la démonstration, encadré et hachuré. Interventions à la main :
« signe » \uncertain{biffé} et remplacé par « "prime" » et « l'effet d'un » ;
« réduits » biffé et remplacé par « canoniques » ; tout l'énoncé (i) « Si
(5.1)(i) est vrai pour $(g,\bar f,\cdot)$ et pour $(g',\pi,\cdot)$ il l'est
pour $(g,f,\cdot)$ » biffé, suivi de « Assertion vide » ; dans (ii), une
ligne tapée biffée et « (5.1(ii)) » repassé ; dans (iii), « et » biffé et une
virgule ajoutée ; le paragraphe « Démonstration. (i) : Supposons d'abord que
le faisceau $F$ sur $X$ est de la forme $F = \pi_*G$. On a » encadré et
hachuré en diagonale}

\page{15}

\[
\mathbf{M}
\]
\[
0 \to G \hookrightarrow G' \twoheadrightarrow G'' \to 0
\]
\[
X_{\mathrm{cris}\,0} \hookrightarrow X_{\mathrm{cris}/n}
\qquad\qquad
\mathcal{M}'' \longrightarrow \mathcal{M}'
\]
\note{les deux dernières expressions sont isolées sur la page, sans texte}

1°) Cristal de prés.\ finie $\mathcal{M}$, sur $S_{\mathrm{cris}}/\mathbb{F}_p$
\struck{\ill{}} avec $F_{\mathcal{M}}$, $V_{\mathcal{M}}$,
\ldots\ d'où

\begin{tikzcd}[column sep=small, row sep=small, nodes={font=\scriptsize}]
 & L^{\prime}(\mathcal{M}) \arrow[dl] & \\
 L_\bullet(\mathcal{M}) \arrow[rr, no head] & & \Delta^*(\mathcal{M})^{(p)} \arrow[ul]
\end{tikzcd}

\qquad $\mathrm{I}(\mathcal{M})$

\begin{tikzcd}[column sep=small, row sep=small, nodes={font=\scriptsize}]
 & L_\bullet(\mathcal{M}) \arrow[dl] & \\
 L'(\mathcal{M}) \arrow[rr, no head] & & \Delta^*(\mathcal{M}) \arrow[ul]
\end{tikzcd}

\qquad $\mathrm{II}(\mathcal{M})$

\note{les flèches horizontales sont de simples traits}

\uncertain{Lui} $\varepsilon^*(\mathcal{M}) = M$

2°) Triangles

\begin{tikzcd}[column sep=small, row sep=small]
 & \mathcal{T}_1 \arrow[dl] & \\
 \mathcal{T}_0 \arrow[rr, no head] & & \mathcal{T}_1 \arrow[ul]
\end{tikzcd}

\quad $\mathcal{T}$ : \qquad (\uncertain{ou} F-V-triangles ?)

\note{le sommet du triangle se lit $\mathcal{T}_1$ comme le coin droit ; l'un
des deux est vraisemblablement un autre indice}

\ill{} Syst.\ de complexes parfaits d'ampl.\ parfaite $\subset [-1,0]$

3°) Isom.\ de (\uncertain{cristaux}) triangles
(\uncertain{ou} F-V-triangles ?)
\[
\varphi^*(\mathcal{T}) \overset{\alpha}{\simeq} \mathrm{I}(\mathcal{M})
\]

4°) Isom.\ (\uncertain{ou} F-V-isom ?)
\[
\mathcal{T}_1 \overset{\beta}{\simeq} \varepsilon^*(\Delta^*(\mathcal{M})) \ \struck{\ill{}}
\]
\uncertain{tel} qu'on ait commutativité

\begin{tikzcd}[column sep=small, row sep=small, nodes={font=\scriptsize}]
\varphi^*(\mathcal{T}_1) \arrow[r, "\alpha_1"] \arrow[dr, "\varphi^*(\beta)"'] & \Delta^*(\mathcal{M})^{(p)} \\
 & \varphi^*\varepsilon^*(\Delta^*(\mathcal{M})) \simeq \Delta^*(\mathcal{M})^{(p)} \arrow[u, Rightarrow, "\mathrm{id}"']
\end{tikzcd}


\page{16}

\note{toute la page est barrée de deux longs traits de crayon en diagonale ;
elle est transcrite sous eux}

\marginal{(3) Indéterminations dans ces données, pour $(\mathcal{M},T)$
(i.e.\ $G$) fixé.}
\note{le « 3 » est cerclé ; un petit tiret orange le précède}

La catégorie des données \struck{\ill{}} en question est « un torseur » sous
la catégorie des torseurs sous $\mu_2$ sur $\mathcal{A}$, donc l'ensemble des
classes d'isomorphisme de ces structures
\struck{\ill{}} \marginal{\uncertain{i.e.\ aussi loc.\ val.\ à l'un des deux-groupes}
$\mathbb{Z}_p^*/\mathbb{Z}_p^{*2}$}
est
\[
H^1(\mathcal{A}, \mu_2) \simeq \prod_p \mathbb{Q}_p^*/\mathbb{Q}_p^{*2} ,
\]
l'ensemble \uncertain{des automorphismes} d'un objet de la catégorie est
\[
H^0(\mathcal{A}, \mu_2) = \prod_p (\mathbb{Z}/2\mathbb{Z}) .
\]
\note{« un torseur sous » est répété sous la ligne et relié par un trait ;
entre les deux formules, une note de marge gauche dont on ne lit que
« \ill{} renforcement \ill{} ! \ill{} »}

(4) Rigidité de $\mathcal{M}$ par les structures envisagées.
\note{le « 4 » est cerclé}

Soit $(\mathcal{M}',T')$ \ill{} des \uncertain{mêmes} données
\uncertain{que} $(\mathcal{M},T)$. On se \uncertain{demande} de déterminer
la catégorie des équivalences
\[
(\mathcal{M}, T, T_{\mathcal{A}'}, u_{\mathcal{A}'}, t_\infty, t_p)
\longrightarrow (\mathcal{M}', T' \ \mathrm{etc.})
\]
\note{au-dessus de $T_{\mathcal{A}'}, u_{\mathcal{A}'}$ une accolade marquée
« $t_{\mathcal{A}'}$ »}
Cette catégorie est vide, ou \struck{elle forme} un torseur sous la catégorie
des torseurs sous $\mu_2$ sur $\mathbb{Q}$, \emph{munis} d'une trivialisation
\struck{\ill{}} au-dessus de $\mathcal{A}$. Cette dernière catégorie a comme
\struck{\ill{}} \ill{} d'objets à isom.\ près le groupe
\[
H^1(\mathbb{Q} \bmod \mathcal{A}, \mu_2) \simeq_2
\Bigl(\prod_p \mathbb{Q}_p\Bigr) \simeq
\Bigl(\prod_p \mathbb{Z}/2\mathbb{Z}\Bigr)/(\pm 1)
\]
\marginal{[qui \uncertain{a le min.\ des conditions}]}
d'autre part cette catégorie est rigide
($H^0(\mathbb{Q} \bmod \mathcal{A}, \mu_2) = 0$).
\note{la lecture « $\mathbb{Q} \bmod \mathcal{A}$ » (cohomologie relative de
$\mathbb{Q}$ modulo les adèles) est douteuse ; le terme du milieu de la
formule se lit « $(\ill{}\ \mathbb{Q}_p)$ »}

Comme l'image de $H^1(\mathbb{Q} \bmod \mathcal{A}, \mu_2)$ dans
$H^1(\mathbb{Q}, \mu_2)$ est nulle, on voit que \emph{l'équivalence
rationnelle entre $(\mathcal{M},T)$ et $(\mathcal{M}',T')$ est déterminée à
isom.\ unique près}. Donc on peut identifier sans ambiguïté $(\mathcal{M},T)$
et $(\mathcal{M}',T')$.

\page{17}

Soit $\omega$ sur $S_{\mathrm{zar}}$, avec $F_\omega$, $V_\omega$, \ldots\
\uncertain{lui associe}

1°) $\mathcal{M} = \varphi^*(\omega)$, \quad
$F_{\mathcal{M}} = \varphi^*(F_\omega)$, \quad
$V_{\mathcal{M}} = \varphi^*(V_\omega)$, \qquad $M = \omega^{p}$,
$\varepsilon^*(\Delta(\mathcal{M})) = \Delta^*(\omega)$
\marginal{complexe can.\ \uncertain{assoc.}\ à $\omega^{\prime}$ par $\omega^{p}$}

2°) Triangles

\begin{tikzcd}[column sep=small, row sep=small, nodes={font=\scriptsize}]
 & {\mathcal{T}_1^{(\omega)} = [\omega \xrightarrow{V_\omega} \omega^{p}]} \arrow[dl, "1"'] & \\
 {\mathcal{T}_0^{(\omega)} = [\omega^{V} \xrightarrow{F_\omega} \omega]} \arrow[rr] & & {\mathcal{T}_1^{(\omega)} = [\omega^{(p)} \xrightarrow{0} \omega^{(p)}]} \arrow[ul]
\end{tikzcd}

\struck{$\mathcal{T}(\omega) \ne$} \quad
$\bigl(\mathcal{T}_1^{(\omega)} = \Delta^*(\omega)\bigr)$ \quad
\uncertain{déduit de $\omega$}
\note{le sommet porte l'indice « 1 » sous une surcharge ; le coin gauche est
lu « $\omega^{V}$ » (peut-être $\omega^{(p)}$)}

3°) Isom.\ \emph{canonique}
\[
\varphi^*(\mathcal{T}(\omega)) \overset{\alpha_\omega}{\simeq}
\mathcal{T}(\varphi^*(\omega))
\]
\[
\mathcal{T}(\omega)_1 \overset{\beta_\omega}{\simeq}
\varepsilon^*\Delta^*(\mathcal{M}) \simeq [\omega \xrightarrow{0} \omega^{p}]
\]

4°) L'\emph{identité}

\emph{Question.} Le foncteur
\[
(\omega, F_\omega, V_\omega) \longmapsto
\bigl(\varphi^*(\omega), \varphi^*(F_\omega), \varphi^*(V_\omega),
\mathcal{T}(\omega), \alpha_\omega, \beta_\omega\bigr)
\]
est-il pl.\ fidèle ?? L'est-il du moins sur la sous-catégorie des triples
$(\omega, F_\omega, V_\omega)$ avec $F_\omega = 0$ ?

\[
\left\lbrace
\begin{array}{l}
\varphi^*(\omega) \xrightarrow{\ u\ } \varphi^*(\omega') \\
\mathcal{T}(\omega) \xrightarrow{\ v\ } \mathcal{T}(\omega')
\end{array}
\right.
\]
\marginal{compatibles avec $F$, $V$ ; hom.\ de \uncertain{triangles}
\uncertain{rendant} commutatifs les diagrammes}
\note{sous $v$, renvoi : « $v_1$ déterminé comme $\varepsilon^*(u)$ » ;
sous le système : « $\varphi^*(v)$ déterminé par $u$ comme $\mathcal{T}(u)$ »}

\begin{tikzcd}[column sep=small, row sep=small, nodes={font=\scriptsize}]
\mathcal{T}_1(\omega) \arrow[r, "v_1"] \arrow[d, Rightarrow, "\beta_\omega"'] & \mathcal{T}_1(\omega') \arrow[d, Rightarrow, "\beta_{\omega'}"] \\
\varepsilon^*\Delta^*(\varphi^*\omega) \arrow[d, no head, "\|" description] & \varepsilon^*(\Delta^*(\varphi^*\omega')) \arrow[d, no head, "\|" description] \\
\Delta^*(\omega) \arrow[r, "\varepsilon^*(u)"'] & \Delta^*(\omega')
\end{tikzcd}

\qquad

\begin{tikzcd}[column sep=small, row sep=small, nodes={font=\scriptsize}]
\varphi^*(\mathcal{T}(\omega)) \arrow[r, no head, "\alpha_\omega \,\sim"] \arrow[d, "\varphi^*(v)"'] & \mathcal{T}(\varphi^*\omega) \arrow[d, "\mathcal{T}(u)"] \\
\varphi^*(\mathcal{T}(\omega')) \arrow[r, no head, "\alpha_{\omega'} \,\sim"'] & \mathcal{T}(\varphi^*\omega')
\end{tikzcd}

\note{les flèches $\beta_\omega$, $\beta_{\omega'}$ sont doubles ; la
première porte en outre un signe surchargé}

\page{19}

\begin{tikzcd}[column sep=small, row sep=small]
 & L' \arrow[dl] & \\
 L_\bullet \arrow[rr] & & \Delta(\mathcal{M}) \arrow[ul]
\end{tikzcd}

\qquad

\begin{tikzcd}[column sep=small, row sep=small, nodes={font=\scriptsize}]
 & \mathcal{M}_{/V} \arrow[dl] & \\
 \mathcal{M}_{/F} \arrow[rr] & & {\mathcal{M}_{/p}^{(\varphi)} = \varphi^*(\Delta(\mathcal{M}))} \arrow[ul]
\end{tikzcd}

\note{sous le premier triangle, une accolade sur $L_\bullet \to \Delta(\mathcal{M})$
marquée « $H_1^{\bullet}$ \quad $H_0$ » ; à droite de $\mathcal{M}_{/V}$, un
mot noirci}

\[
0 \to n \to \mathbf{D}_1 \to t^* \to \omega \to \mathbf{D}_0 \to \nu^* \to 0
\]
\[
0 \to {}_F\tilde{\mathcal{M}} \to {}_{F}\tilde{\mathcal{M}} \xrightarrow{F}
{}_V\mathcal{M} \to \mathcal{M}_F \xrightarrow{V} \tilde{\mathcal{M}}_p \to
\tilde{\mathcal{M}}_V \to 0
\]
\note{la seconde suite est lue par morceaux : au-dessus de ${}_V\mathcal{M} \to
\mathcal{M}_F$ une flèche coudée passe par un $\Theta$ (en crayon) ; les
indices et les tildes sont incertains}

\[
X_{\mathrm{cris}} \ \downarrow \varphi \ X_{\mathrm{zar}}
\]

\struck{(a) descendre $\mathcal{M}_{/F}$ et $\mathcal{M}_{/p}$}

\struck{(b) descendre le triangle tout entier}
\note{les deux lignes sont biffées d'un trait en zigzag}

\begin{tikzcd}[column sep=small, row sep=small, nodes={font=\scriptsize}]
J \arrow[r, hook] \arrow[d] & B \arrow[r] \arrow[d, "\mu"] & A \arrow[d] \arrow[dr, "{\lambda \mapsto \lambda^p}"] & \\
JB' \arrow[r, hook] & B' \arrow[r] & B'/JB' \arrow[r] \arrow[ddr] & A \\
J^{(p)} & B \arrow[d, dashed] & B/J^{(p)}B \arrow[d] & \\
 & {} & B \arrow[r] & C \arrow[uu]
\end{tikzcd}

\note{diagramme très retouché : sous $JB'$, $B'$, $B'/JB'$, des signes
d'égalité verticaux mènent à $J^{(p)}$, $B$, $B/J^{(p)}B$ ; « $\gamma^p$ »
devant $JB'$ ; une flèche en pointillé courbe vers le haut près de $B$ ; une
colonne sous $B$ est noircie ; une flèche biffée $C \to C$ ; l'arête courbe
de $A$ vers $A$ est hachurée ; la case vide du tikzcd tient la place du
pointillé}

\[
B - A, \qquad B \otimes B \to A
\]

\begin{tikzcd}[column sep=small, row sep=small, nodes={font=\scriptsize}]
J^{(p)}B \arrow[dr] & J \arrow[r, hook] \arrow[dr] & B \arrow[r, no head] \arrow[dr] & A \\
 & 0 & I \arrow[r, hook] & C' \arrow[r] & A
\end{tikzcd}

\quad « p.d. » sous $I$

\begin{tikzcd}[column sep=small, row sep=small, nodes={font=\scriptsize}]
 & & \Delta^*(\mathcal{M}) \arrow[dl, no head] \arrow[dr, no head] & \\
 & \tilde{D}_1 & & \tilde{D}_0
\end{tikzcd}

\[
0 \to n \to D_1 \to \tilde{D}_1 \to 0, \qquad
0 \to \tilde{D}_0 \to D_0 \to \nu^* \to 0
\]
\marginal{[uniques \uncertain{déterminés} par $\mathrm{I}(\mathcal{M})$]}

\[
(\mathrm{i}) \qquad 0 \to \tilde{D}_1 \to t^* \to \omega \to \tilde{D}_0 \to 0
\]
\note{au-dessus, $t^* \to \vartheta \to \omega$ par une flèche coudée passant
par $\vartheta$ ; au-dessus de $\tilde D_1$ un mot noirci}

\marginal{On descend $\Theta$ en $\vartheta$, et \struck{\ill{}} \ill{}
\ill{} en extension de $\tilde D'_0$ par $\vartheta$ \struck{\ill{}} et de
$\vartheta$ par $D'_1$}
\note{cette note est séparée des diagrammes par un trait vertical ondulé ;
en dessous, un « A » noirci}

(ii)

\begin{tikzcd}[column sep=small, row sep=small]
 & t^*[1] \arrow[dl, "1"'] & \\
 \omega^* \arrow[rr] & & \tilde{\Delta}^*(\mathcal{M}) \arrow[ul]
\end{tikzcd}

\note{une flèche courbe marquée
« connu » y renvoie ; un signe noirci à sa droite}

\marginal{\struck{\ill{} explic\ill{}} descend : \quad
$\mathcal{M}_F \to \varphi^*(\tilde\Delta^*(\mathcal{M}))$, \quad
$\mathcal{M}_{/V}[1]$, \quad (« $\mathcal{M}_{/p}$ »)$'$}

i.e.\ on doit \struck{donner} \uncertain{considérer} les isom.
\[
[t^* \to \omega]\,\struck{\ill{}} \simeq \tilde{\Delta}^*(\mathcal{M})
\]
descendant \struck{\ill{}} \uncertain{ceux} \ill{} par $\varphi^*$ de
$\Delta^*(\mathcal{M})$ \ill{} \uncertain{isom.}\ \ill{} à \ill{} près, ou
\ill{} $\varphi^*$ dans

l'arbitraire est dans \ill{} \ill{}, \ill{} principalement, sans doute,
i.e.\ dans $\operatorname{Ext}^1(\tilde D_0, \tilde D_1)$ ou plus
\ill{} dans
\[
\operatorname{Ext}^1(S_{\mathrm{zar}} \bmod S_{\mathrm{cris}};
\tilde D_0, \tilde D_1) ;
\]
les \uncertain{autres} de la structure \ill{} \ill{}
\[
\operatorname{Ext}^0(S_{\mathrm{zar}} \bmod S_{\mathrm{cris}};
\tilde D_0, \tilde D_1) = 0 \quad !
\]
\marginal{Obstruction à \uncertain{changer} \ldots}
dans
\[
\operatorname{Ext}^2(S_{\mathrm{zar}} \bmod S_{\mathrm{cris}},
\tilde D_0, \tilde D_1) .
\]

\note{en haut à droite de la moitié inférieure, un diagramme de schémas :}
\[
\begin{array}{ccccc}
\tilde{X} & \equiv & \widetilde{X \times X} & \equiv & \ldots \\
| & & & & \\
X & = & X \times X & \equiv & X \times X \times X \\
| & & & & \\
S' & = & S & = & S \\
\downarrow & & \downarrow & & \downarrow \\
\tilde{X} & \Leftarrow & \widetilde{X \times X} & \Lleftarrow & \widetilde{X \times X \times X} \\
\downarrow & & \downarrow & & \downarrow \\
\tilde{X} & \Leftarrow & X \times X & \Lleftarrow & X \times X \times X \\
\downarrow \mathrm{plat} & & \downarrow \mathrm{plat} & & \downarrow \mathrm{plat} \\
S & = & S & = & S
\end{array}
\]
\note{la lecture est approchée : à gauche « $S \hookleftarrow X$ »,
$\hookleftarrow$ entre $\tilde X$ et $X$ ; des flèches en pointillé
bouclent sur chaque colonne ; une grande flèche courbe marquée « $F_S$ »
descend de $\tilde X$ jusqu'au $S$ du bas ; les tildes larges couvrent
chaque produit}

\end{document}
